\section{Corrections to the record}
\label{sec:corrections}

This project's convention is that corrections are stated as corrections, in their own section, and not
silently patched into the artifacts they correct. The \vsix{} study opened a register against the
\vfour{} and \vfive{} record; this cycle adds the entries below. Each names the measurement that forces
it.

\subsection{Corrections to the v0.4--v0.6 record}
\label{sec:corrections-corpus}

\begin{description}

\item[The cost of the frontier's search, wrong in both directions.] The corpus recorded the searching configuration at \vsevenCorpusCost$\times$ the
  blueprint, and the \vsix{} paper at ``three orders of magnitude''. Both figures describe the
  \emph{unrestricted} configuration, which is not the point of the frontier anything was measured at.
  The endgame-restricted operating point was re-measured at $\sim$\vsevenCostRecorded$\times$ during
  this cycle, and that correction was carried into the protocol. \textbf{This report corrects the correction}: measured on the shared panel at a fixed thread count,
  \fcheap{} --- the configuration the \vsevenCostRecorded$\times$ describes --- costs at least
  \vsevenCostFcheapTight$\times$ the blueprint, and \frozen{} at least
  \vsevenCostFrozenTight$\times$ (\S\ref{sec:results-cost}). Every figure in that table is a lower
  bound and none carries an interval. The record has now been wrong in both directions and the reason is the
  same both times --- a ratio quoted without the basis it was computed on.

\item[The corpus's exploitability figure is a lower bound, and has been read as more.] The v0.4 findings document reports that a policy fitted
  specifically against v0.4 fails to reach parity with it; the same study's own table records the
  response reaching $\vsevenVfourLbr\%$ with an interval spanning parity. The \vsix{} study corrected this. It is
  repeated here because this cycle's own strongest result --- \S\ref{sec:exploit}, eight searches, all
  losing --- is a number of exactly the same kind, and the correction is what stops this paper making
  the same claim.

\item[The declaration-allocation channel is worth zero, not two points.] The
  register ranked it second by priority, at \vsevenLonePriority{}. \S\ref{sec:rejected-alloc} closes it
  at exactly zero over $\sim$\vsevenLoneDecls{} recorded declarations, shows its proposed fix is worse
  than the object it would replace, and shows most of the residual error class is an information limit
  rather than a defect.

\item[The forced-endgame entry's own number was wrong.]
  Re-derived from this cycle's capture, closing the entire accuracy gap is worth
  \vsevenLthirteenClose{}~win-rate points --- the register's figure corrected \emph{upward} by a factor
  of two, and still a fiftieth of the detection floor. The entry is closed either way.

\item[The detection floor does not buy down with evaluation games.] Phase~1 measured
  \vsevenFloorPhaseOne{} and predicted the floor would fall as the inverse square root of evaluation
  games. At exactly four times the games it is \vsevenPrimaryFloor{}, and the sub-floor rung is still
  undetected because the excess resolves \emph{to} zero rather than because the interval is wide. Below
  roughly a point and a half the responder is empty, not unresolved. This report reproduces that on
  holdout material (\S\ref{sec:results-controls}).

\end{description}

\subsection{Corrections this cycle makes to its own record}
\label{sec:corrections-self}

\begin{description}

\item[A freeze-verification step could not fire, and its silence was not a pass.] The freeze checker reads its source-drift baseline out of the freeze artifact's own
  provenance field, and that artifact was rewritten twice after the freeze commit, each rewrite
  refreshing the baseline to the then-current tree. The check therefore compared the tree against a
  baseline taken from the tree, and the sources the protocol predicted would be listed as changed no
  longer differed from it. Measured against the actual freeze commit, \vsevenDriftChanged{} of the
  \vsevenDriftTotal{} sources differ. \textbf{What is not weakened is the policy identity}: the round-trip assertions and the
  mirror digest are computed by \emph{running} the engine, not by comparing file hashes, and all of
  them round-trip. Because the hash-based check was vacuous it was replaced by an executed one --- the
  engine was rebuilt from the tree as it stood at the end of the battery and the mirror digest
  recomputed with the fresh binary. It is \texttt{\vsevenMirrorDigest{}}, identical to the frozen
  artifact and to the binary that played every cell.

\item[A scored cell is not bit-reproducible for a searching arm.]
  \label{sec:corrections-reproducibility}
  Before any holdout material was spent, one training cell was checked to be bit-identical over
  \vsevenReproRuns{} runs across three thread counts, and it was
  (\vsevenReproIdentical{}). \textbf{That check was too small.} The panel contains three policies
  twice, so twenty-four pairs of byte-identical commands were run twice each; one pair disagreed, by
  one game in \vsevenResidueGames{}. That is the cross-deal agent residue of \S\ref{sec:results-residual}
  reaching the \emph{scored} mode, which four hundred deals was not large enough to expose. The
  magnitude is roughly \vsevenResidueRatio$\times$ below a cell half-width and no verdict in this paper turns on
  it, but the earlier claim should be read as ``reproducible at four hundred deals'', not
  ``reproducible''.

\item[The scored-cell count is not the count the artifacts support.]
  That document states \vsevenScoredCellsRecorded{} scored match cells in three places; the figure is
  hand-typed in the reduction script rather than computed. Counting rows that carry a match object
  across every scored battery gives \textbf{\vsevenScoredCells{}}. Nothing turns on it --- it is the
  denominator of an observation about a single action-limit hit --- but a corpus whose standing rule is
  that no number is hand-typed should not carry a hand-typed total, and the generator for this paper
  now computes it.

\item[The partner-regime table was published without its own intervals.]
  \label{sec:corrections-intervals}
  The phase-5 results document prints the partner deltas as bare point estimates with a
  replication flag. Every underlying cell carries a deal-clustered interval, and the delta arithmetic
  is fixed by the protocol, so the intervals are computable and were simply compressed out of a
  markdown table. Table~\ref{tab:partners} restores them. No verdict changes; a reader's ability to
  judge the verdict does.

\item[The cross-play gap was printed under the reversed sign convention.] The phase-5 document reports the gap as \vsevenXpGapRecorded{}, which is diagonal minus
  off-diagonal. The claim it settles asks whether the \emph{off}-diagonal collapses, so the informative direction is off-diagonal minus diagonal, which is \vsevenXpGap{}.
  \emph{Neither sign is resolved}: the magnitude is \vsevenXpGapAbs{}~points against a per-cell
  half-width of \vsevenXpHalfWidth{} and a \vsevenStwoThreshold{}-point threshold, so no verdict
  moves and this paper reads no lift out of it either. What the correction is worth is only this:
  the sign as printed \emph{implies} a small collapse, and the data support no direction at all.

\end{description}

\subsection{A note on comparability}
\label{sec:corrections-comparability}

Two things about this cycle's evidence standard are not the corpus's, and a reader comparing across
versions should have them. Every claim here is measured on material sealed before the configuration
existed, under a protocol committed before any of it was played, with a stated failure condition ---
which no previous study in this corpus did. And every claim is replicated across two disjoint banks
with the sign rule of \S\ref{sec:methodology-intervals} applied, which is why
\S\ref{sec:against-location} reports a non-replication that a pooled-interval convention would have
reported as a pass.
